\documentclass[11pt]{article}

\usepackage[margin=1in]{geometry}
\usepackage{amsmath}
\usepackage{graphicx}
\usepackage{float}

\setlength{\parindent}{0pt}
\setlength{\parskip}{6pt}

\title{\textbf{COMP SCI 640 --- Lab 4 Report}}
\author{}
\date{}

% Compile with: pdflatex report_lab4.tex
% Replace the blank spaces with your answers; add pages as needed.
% The comments below are reminders. Follow Q1--Q5 in the assignment specification.
% Cite run IDs, log filenames, and line numbers for your evidence.
% Include selected annotated excerpts in this report; submit full logs separately.

\begin{document}
\maketitle

\section*{Group Information}

\textbf{Group Number:} \underline{\hspace{3cm}}

\textbf{Group Members:}
\begin{enumerate}
    \item \underline{\hspace{7cm}}
    \item \underline{\hspace{7cm}}
    \item \underline{\hspace{7cm}}
    \item \underline{\hspace{7cm}}
\end{enumerate}

\section*{Setup Photo}

% Put a photo of both boards in the same folder and name it boards.jpg.
% It will replace the placeholder below automatically.
% Change the filename or image dimensions here if needed.
\begin{figure}[H]
    \centering
    \IfFileExists{boards.jpg}{%
        \includegraphics[width=0.9\linewidth,height=5cm,keepaspectratio]{boards.jpg}%
    }{%
        \fbox{\parbox[c][4cm][c]{0.9\linewidth}{\centering
            Insert a photo showing both boards.}}%
    }
    \caption{Our two-board experimental setup.}
\end{figure}

\newpage

\subsection*{Q1}
% Write your answer here.
% Decode one actual 20-byte DATA header and show successful CRC verification.
% Show application delivery and verification of the final 17-byte message.
% With a two-second initial send delay, show the application waiting and waking.
% Cite the blocking queue-receive call in your code.
\vspace{2cm}

\subsection*{Q2}
% Write your answer here.
% Annotate the normal-delivery trace to show several outstanding DATA packets.
% For the DATA 3 loss, show repeated ACKs, discarded out-of-order packets,
% the timeout, and the retransmitted range.
% For the corrupted final ACK, show received/recomputed CRCs, rejection before
% state changes, and recovery without duplicate application delivery.
% Compare the observed timeout with 500 ms using one board's timestamps.
\vspace{2cm}

\subsection*{Q3}
% Write your answer here and insert your queue/window plot.
% Plot queue occupancy Q and advertised Window W using the receiver's clock.
% Identify a packet acknowledged before the application read it; mark both events.
% Show one sender-side ACK (a,w), the resulting limit R, and newly allowed DATA.
% Explain overflow prevention and compare the fast and slow runs.
% State whether you used an initial pause in application reads.
%
% Example for inserting your plot (uncomment after adding the image):
% \begin{figure}[H]
%     \centering
%     \includegraphics[width=\linewidth]{queue_window.pdf}
%     \caption{Queue occupancy and advertised window during our slow run.}
%     \label{fig:queue-window}
% \end{figure}
\vspace{2cm}

\subsection*{Q4}
% Write your answer here.
% Show ACK (8,0), one application read, the dropped ACK (8,1), the periodic
% replacement, and acceptance of DATA 9. Explain why sending resumes even
% though the ACK number stays at 8.
% Show transfer completion and that no second application read broke the stall.
\vspace{2cm}

\subsection*{Q5}
% Write your answer here and include the F/S/L comparison table.
% For each condition, report: successful trials out of three; median and
% minimum--maximum successful goodput (bytes/s); median DATA retransmissions
% across all three trials; and the largest queue occupancy observed.
% Use N/A for goodput when no trial succeeded.
% Explain performance using your measurements and earlier traces. Include
% integrity checks and full-queue discard counts.
% Describe one bug fixed or a specific correctness check, with evidence.
% Submit all nine trial rows separately in the required CSV file.
\vspace{2cm}

\end{document}
